\section{RESTful Services}

\begin{frame}{Section 6.3: RESTful Services}
\begin{itemize}
    \item Microservices may communicate synchronously or asynchronously using different protocols.
    \item One of the most widely used approaches for service interaction is \textbf{REST}.
    \item RESTful services follow a set of conventions built on top of HTTP.
    \item In practice, REST is one of the most common ways for microservices to exchange information.
\end{itemize}
\end{frame}

\begin{frame}{What REST Really Is}
\begin{itemize}
    \item Strictly speaking, REST is not a single formal protocol.
    \item It is an \textbf{architectural style} for service communication.
    \item RESTful services:
    \begin{itemize}
        \item use HTTP,
        \item organize information as resources,
        \item exchange representations of those resources.
    \end{itemize}
    \item REST stands for \textbf{Representational State Transfer}.
\end{itemize}
\end{frame}

\begin{frame}{Core Idea Behind REST}
\begin{itemize}
    \item REST is based on transferring representations of digital resources between server and client.
    \item A resource could be:
    \begin{itemize}
        \item a book,
        \item a customer record,
        \item a traffic incident,
        \item an uploaded image,
        \item an order.
    \end{itemize}
    \item The same resource may have different representations, such as JSON, XML, or HTML.
\end{itemize}
\end{frame}

\begin{frame}{RESTful Service Principles}
\begin{itemize}
    \item \textbf{Use HTTP verbs} to express operations.
    \item \textbf{Be stateless:} the server should not rely on remembered client session state.
    \item \textbf{Use URI-addressable resources:} each resource should have a unique identifier.
    \item \textbf{Use structured representations:} typically JSON, and sometimes XML.
\end{itemize}
\end{frame}

\begin{frame}{Why Statelessness Matters}
\begin{itemize}
    \item In a stateless service, each request contains the information needed to process it.
    \item The server does not depend on stored interaction state from earlier requests.
    \item This supports:
    \begin{itemize}
        \item scalability,
        \item easier replication,
        \item simpler failure recovery.
    \end{itemize}
    \item Statelessness fits well with cloud-based microservices.
\end{itemize}
\end{frame}

\begin{frame}{Resources and URIs}
\begin{itemize}
    \item In REST, a resource is accessed through a \textbf{URI}.
    \item URIs should reflect a meaningful hierarchical structure.
    \item Example idea:
    \begin{itemize}
        \item road incidents organized by road, place, direction, and incident number
    \end{itemize}
    \item A good URI structure makes resources easier to identify and operate on.
\end{itemize}
\end{frame}

\begin{frame}{CRUD Operations in REST}
\begin{itemize}
    \item RESTful services commonly support four fundamental operations:
    \begin{itemize}
        \item \textbf{Create}
        \item \textbf{Read}
        \item \textbf{Update}
        \item \textbf{Delete}
    \end{itemize}
    \item These are often mapped onto standard HTTP verbs.
\end{itemize}
\end{frame}

\begin{frame}{HTTP Verbs and REST Operations}
\begin{table}
\centering
\begin{tabular}{ll}
\toprule
Operation & Typical HTTP Verb \\
\midrule
Create & POST \\
Read & GET \\
Update & PUT \\
Delete & DELETE \\
\bottomrule
\end{tabular}
\end{table}

\vspace{0.5em}
\begin{itemize}
    \item This mapping gives RESTful services a consistent structure.
\end{itemize}
\end{frame}

\begin{frame}{Example: Traffic Incident Service}
\begin{itemize}
    \item Imagine a service that stores traffic incidents.
    \item Example operations might include:
    \begin{itemize}
        \item \textbf{GET} to retrieve incident information,
        \item \textbf{POST} to add a new incident,
        \item \textbf{PUT} to update an incident,
        \item \textbf{DELETE} to remove an incident.
    \end{itemize}
    \item This is a good example of a resource-oriented REST design.
\end{itemize}
\end{frame}

\begin{frame}{Request and Response Processing}
\begin{itemize}
    \item A RESTful microservice accepts an HTTP request,
    \item processes the request,
    \item performs the required service actions,
    \item and generates an HTTP response.
    \item This creates a predictable request/response flow.
\end{itemize}
\end{frame}

\begin{frame}{HTTP Request and Response Structure}
\begin{itemize}
    \item A typical HTTP request includes:
    \begin{itemize}
        \item an HTTP verb,
        \item a URI,
        \item an HTTP version,
        \item request headers,
        \item optionally, a request body.
    \end{itemize}
    \item A typical HTTP response includes:
    \begin{itemize}
        \item an HTTP version,
        \item a response code,
        \item response headers,
        \item optionally, a response body.
    \end{itemize}
\end{itemize}
\end{frame}

\begin{frame}{Important HTTP Header Information}
\begin{itemize}
    \item \textbf{Accept} specifies what response content types the client can handle.
    \item \textbf{Content-Type} specifies the type of the request or response body.
    \item \textbf{Content-Length} specifies the size of the message body.
    \item These headers help clients and services interpret exchanged data correctly.
\end{itemize}
\end{frame}

\begin{frame}{JSON vs. XML}
\begin{itemize}
    \item RESTful services often return structured data in \textbf{JSON}.
    \item XML can also be used, but JSON is usually preferred in modern systems.
    \item JSON is popular because it is:
    \begin{itemize}
        \item simpler,
        \item lighter weight,
        \item easier to parse in many environments.
    \end{itemize}
    \item XML is more flexible, but often more verbose.
\end{itemize}
\end{frame}

\begin{frame}{Example Resource Representation}
\begin{itemize}
    \item A traffic incident resource might include:
    \begin{itemize}
        \item incident ID,
        \item date,
        \item time,
        \item road ID,
        \item place,
        \item direction,
        \item severity,
        \item description.
    \end{itemize}
    \item REST returns a representation of that resource, not just a vague blob of text.
\end{itemize}
\end{frame}

\begin{frame}{Why REST Works Well for Microservices}
\begin{itemize}
    \item REST is widely understood and widely supported.
    \item It fits naturally with web technologies and HTTP infrastructure.
    \item It encourages:
    \begin{itemize}
        \item simple interfaces,
        \item stateless interaction,
        \item consistent operations on resources.
    \end{itemize}
    \item This makes it a practical communication style for many microservice systems.
\end{itemize}
\end{frame}

\begin{frame}{Limits of REST}
\begin{itemize}
    \item REST is useful, but it is not always the only or best choice.
    \item Some systems may need:
    \begin{itemize}
        \item asynchronous messaging,
        \item event-based communication,
        \item specialized protocols.
    \end{itemize}
    \item REST is best seen as one important tool in the microservices toolbox.
\end{itemize}
\end{frame}

\begin{frame}{Section 6.3 Recap}
\begin{itemize}
    \item RESTful services use HTTP and treat information as resources.
    \item Resources are identified by URIs and manipulated using standard HTTP verbs.
    \item REST emphasizes stateless interaction and structured representations.
    \item JSON is commonly preferred for modern microservice communication.
\end{itemize}
\end{frame}

\begin{frame}{Looking Ahead}
\centering
\Large Next: Section 6.4 Service Deployment

\vspace{1em}

\normalsize
We next examine what happens after development, when services must actually be deployed, updated, and operated.
\end{frame}
